\section{Introduction}
\label{sec:intro}

Multi-teacher on-policy distillation (MOPD) transfers domain specialists
into one student by scoring student-generated responses with their
assigned teachers \citep{ma2026mopdmultiteacheronpolicydistillation}.
The objective is to integrate their capabilities, but equal numbers
of prompts from each domain need not give those domains equal influence
on training. Open-MOPD identifies capability imbalance and studies
token-share weighting, changing distillation gaps, and stale
student-dependent rewards \citep{gao2026openmopddiagnosingfixingcapability}.
These findings motivate examining how the objective weights each
domain and which student parameters its teacher changes most.

Prior work reports sparse cumulative parameter changes in RL
\citep{mukherjee2025subnetworks} and OPD
\citep{yu2026denseupdates,shen2026opdgeometry}.
We use update sparsity to describe concentration of actual parameter
changes, distinguishing individual optimizer steps from cumulative
changes over training. We ask how multiple teachers affect a common
student during MOPD, and use gradients as an additional measurement to
separate the current loss signal from optimizer history. Independently
trained students provide a reference, but their endpoints do not identify
teacher contributions within the joint student.

After introducing the MOPD learning objective in
Section~\ref{sec:setting}, we develop three studies and place each
experiment alongside its corresponding question:
\begin{enumerate}
\item \textbf{Token balancing (Section~\ref{sec:normalization}).}
We derive how global-token, domain-token, and domain-response averaging
weight domains and response lengths, then compare their effects on
capability balance. The analysis explains existing loss reductions
without introducing a new allocation algorithm.
\item \textbf{Update sparsity (Section~\ref{sec:subnetworks}).}
We measure the concentration of optimizer steps and cumulative parameter
changes in single-teacher and joint OPD, compare the parameters most
changed by each teacher at the same MOPD checkpoint, and relate their
overlap to teacher distribution differences. Raw gradients help
interpret these updates.
\item \textbf{Supervision density (Section~\ref{sec:density}).}
We compare sampled-token PG and teacher top-$k$ training in both
settings, with full-vocabulary updates on common prefixes as a
local reference. We test whether update sparsity differs materially
and report capability alongside it.
\end{enumerate}

Shared parameter locations can receive compatible or opposing changes,
and sampled-token feedback still depends on all logits through the
softmax normalizer. Neither low overlap nor limited vocabulary coverage
therefore guarantees sparse parameter updates or better integration.
\missing{Replace the prospective empirical contributions with observed
findings and their measured scope once the experiments are complete.}
